\documentclass[review,1p,times]{elsarticle}
\usepackage{booktabs}
\usepackage{multirow}
\usepackage{longtable}
\usepackage{array}
\usepackage{siunitx}
\usepackage{amsmath,amssymb}
\usepackage[hidelinks]{hyperref}
\newcolumntype{L}[1]{>{\raggedright\arraybackslash}p{#1}}
\begin{document}
\section*{Supplementary Material: Migrated Supporting Tables}
This file contains non-core tables migrated from the main manuscript to enforce the four main-matrix structure in the primary text.

% Supplementary table 1: tab:pimoe_compare
\begin{table}[t]
\centering
\caption{Comparison of the baseline Ultra-LSNT and the FG-MoE feasibility-clipper extension.}
\label{tab:pimoe_compare}
\begingroup
\setlength{\tabcolsep}{4pt}
\small
\resizebox{\linewidth}{!}{%
\begin{tabular}{L{0.16\linewidth}L{0.40\linewidth}L{0.40\linewidth}}
\toprule
Dimension & Current Ultra-LSNT & FG-MoE feasibility clipper (upgrade path) \\
\midrule
Interpretability source & Post-hoc correlation between routing/expert usage and covariates (audit signal) & Empirical envelope heuristic with plausibility penalties; not causal mechanism discovery \\
Signal decomposition & Moving-average (MA) decomposition for trend/season separation & Moving-average (MA) decomposition (unchanged); optional VMD/EEMD extensions are discussed in the main manuscript discussion section. \\
Multi-scale encoding & Data-driven Conv1D kernels with multi-scale temporal embedding & Same multi-scale encoder; with an additional output plausibility barrier via $P_{\mathrm{env}}(\hat v)$ \\
Robustness mechanism & Soft sparse routing (noise filtering via conditional computation) & Soft routing plus static envelope clipping (power-envelope barrier and inference projection) \\
\bottomrule
\end{tabular}%
}
\endgroup
\end{table}

% Supplementary table 2: tab:dispatch_params
\begin{table}[t]
\centering
\caption{Dispatch model parameters for the 120-day rolling day-ahead evaluation on Wind (CN). Power is in kW, energy in kWh, and costs in CNY.}
\label{tab:dispatch_params}
\begin{tabular}{ll}
\toprule
Parameter & Value \\
\midrule
Time resolution / horizon & $\Delta t=1$\,h; $T=24$ (rolling evaluation) \\
Demand profile $d_t$ & Fixed demand of 50{,}000\,kW (constant across $t$) \\
Wind uncertainty rule & UQ-aware adaptive: $\sigma_{w,t}=\sigma_{\mathrm{base}}+\alpha_H H_t + \alpha_D D_t$ (kW), with fallback to fixed $\sigma_{w,t}=0.2|\hat{w}_t|$ for baseline comparability \\
Thermal generation bounds & $g_{\min}=10{,}000$\,kW; $g_{\max}=80{,}000$\,kW \\
Thermal ramp constraint & $R_{\max}=20{,}000$\,kW/h \\
Fuel cost & 0.5 CNY/kWh \\
Startup cost coefficient $c_{\text{su}}$ & 5{,}000 CNY per start \\
Shutdown cost coefficient & Same as $c_{\text{su}}$ (implementation uses startup\_cost) \\
Reserve cost & 0.2 CNY/kWh (reserve capacity) \\
Storage energy capacity & $S_{\max}=50{,}000$\,kWh \\
Storage power limits & $P_{\text{ch},\max}=10{,}000$\,kW; $P_{\text{dis},\max}=10{,}000$\,kW \\
Storage efficiencies & $\eta_{\text{ch}}=0.95$; $\eta_{\text{dis}}=0.95$ \\
Storage operating cost & 0.05 CNY/kWh \\
Initial state of charge & $\text{SOC}_0=0.5$ (fraction of $S_{\max}$) \\
\bottomrule
\end{tabular}
\end{table}

% Supplementary table 3: tab:dispatch_mapping_sensitivity
\begin{table}[t]
\centering
\caption{Dispatch mapping sensitivity from archived RTS-24 DC-OPF aggregates (single added rerun at $w_{\max}=300$ MW). Results compare three mapped corridors under the same fixed mapped-and-clipped screening protocol.}
\label{tab:dispatch_mapping_sensitivity}
\begingroup
\setlength{\tabcolsep}{4pt}
\small
\begin{tabular}{llccc}
\toprule
Mapping case & Model & Total RT cost (M CNY) & Curtailment (MWh) & Curtailment gain vs DLinear (\%) \\
\midrule
wmax260 & DLinear & 137.08 & 61801.35 & +0.00 \\
wmax260 & iTransformer & 145.43 & 54242.21 & +12.23 \\
wmax260 & TimeMixer & 147.91 & 54242.21 & +12.23 \\
wmax260 & Ultra-LSNT & 145.59 & 54381.53 & +12.01 \\
\midrule
wmax280 & DLinear & 133.14 & 97778.62 & +0.00 \\
wmax280 & iTransformer & 134.50 & 93914.96 & +3.95 \\
wmax280 & TimeMixer & 134.69 & 93914.96 & +3.95 \\
wmax280 & Ultra-LSNT & 134.61 & 93930.40 & +3.94 \\
\midrule
wmax300 & DLinear & 130.07 & 139403.69 & +0.00 \\
wmax300 & iTransformer & 130.06 & 138093.39 & +0.94 \\
wmax300 & TimeMixer & 130.06 & 138093.39 & +0.94 \\
wmax300 & Ultra-LSNT & 130.06 & 138093.39 & +0.94 \\
\bottomrule
\end{tabular}
\vspace{3pt}
\footnotesize Source: \path{results/supplementary_evidence/dispatch_mapping_sensitivity_wmax260_280_300.csv}. This is mapped-interval sensitivity evidence only; all conclusions remain conditional on the chosen mapped-and-clipped screening interface.
\endgroup
\end{table}

% Supplementary table 3b: tab:dispatch_mapping_daylevel_consistency
\begin{table}[t]
\centering
\caption{Dispatch mapping day-level consistency under three mapped corridors ($w_{\max}=260/280/300$ MW). Values report daily win frequency versus DLinear and median daily deltas under the same fixed screening interface family.}
\label{tab:dispatch_mapping_daylevel_consistency}
\begingroup
\setlength{\tabcolsep}{4pt}
\small
\begin{tabular}{llccccc}
\toprule
Mapping case & Model & Cost day-win vs DLinear (\%) & Curtailment day-win vs DLinear (\%) & Median $\Delta$Cost (M CNY/day) & Median $\Delta$Curtailment (MWh/day) & $n_{\text{days}}$ \\
\midrule
wmax260 & iTransformer & 41.67 & 83.33 & +0.0012 & -0.2418 & 120 \\
wmax260 & TimeMixer & 35.83 & 71.67 & +0.0012 & -0.2418 & 120 \\
wmax260 & Ultra-LSNT & 42.50 & 85.83 & +0.0012 & -0.2418 & 120 \\
\midrule
wmax280 & iTransformer & 30.83 & 48.33 & +0.0000 & +0.0000 & 120 \\
wmax280 & TimeMixer & 49.17 & 71.67 & +0.0000 & +0.0000 & 120 \\
wmax280 & Ultra-LSNT & 35.00 & 48.33 & +0.0000 & +0.0000 & 120 \\
\midrule
wmax300 & iTransformer & 61.67 & 57.50 & +0.0000 & +0.0000 & 120 \\
wmax300 & TimeMixer & 54.17 & 50.00 & +0.0000 & +0.0000 & 120 \\
wmax300 & Ultra-LSNT & 58.33 & 51.67 & +0.0000 & +0.0000 & 120 \\
\bottomrule
\end{tabular}
\vspace{3pt}
\footnotesize Source: \path{results/supplementary_evidence/dispatch_mapping_daylevel_consistency_wmax260_280_300.csv}. Day-win captures daily relative advantage frequency, and aggregate directional gains can diverge from day-level stability. Evidence remains conditional on the mapped-and-clipped screening protocol and is not a direct operator-readiness proof.
\endgroup
\end{table}

% Supplementary table 3c: tab:supp_decision_aware_robustness_summary
\begin{table}[t]
\centering
\caption{Consolidated cross-regime metrics for four principal models (supplementary screening view).}
\label{tab:supp_decision_aware_robustness_summary}
\begingroup
\setlength{\tabcolsep}{3pt}
\scriptsize
\begin{tabular*}{\linewidth}{@{\extracolsep{\fill}}lcccccccc@{}}
\toprule
\textbf{Model} & \textbf{$R^2@0.60$} & \textbf{Miss. $R^2$} & \textbf{L1 FB} & \textbf{L1 Cost} & \textbf{L2 Curt.} & \textbf{L2 Succ.} & \textbf{Rank A/V/W} & \textbf{Backfire} \\
\midrule
DLinear & 0.3471 & 0.9245 & 62.5 & 54.47 & 139404 & 100.0 & 4.00/0.00/4 & 0.00 \\
TimeMixer & 0.8417 & 0.8731 & 21.7 & 34.90 & 138093 & 100.0 & 2.33/0.22/3 & 49.70 \\
iTransformer & 0.8580 & 0.8777 & 16.7 & 32.16 & 138093 & 100.0 & 1.33/0.22/2 & 46.10 \\
Ultra-LSNT & 0.8180 & 0.8446 & 24.2 & 35.97 & 138093 & 100.0 & 2.33/0.89/3 & 45.31 \\
\bottomrule
\end{tabular*}
\vspace{3pt}
\footnotesize Source mapping: $R^2@0.60$ and Missingness from \path{results/supplementary_evidence/fairness_closure_fullfault_hard_table.csv} with \texttt{training\_mode=fault\_augmented}; L1 fallback from \path{results/supplementary_evidence/dispatch_raw_vs_mapped_summary.csv}; L1 cost from \path{results/supplementary_evidence/main4_matrix4_dispatch_value.csv} (\texttt{layer=L1\_CopperPlate\_120d}); L2 curtailment/success from \path{results/supplementary_evidence/dispatch_closure_mapping_decision_aggregate.csv} (\texttt{mapping=affine, subset=all}); rank robustness from \path{results/supplementary_evidence/dispatch_closure_mapping_decision_ranking_stability.csv}; Backfire(\%) is mean \texttt{backfire\_rate\_vs\_dlinear} across three mappings in \texttt{subset=all}. Consolidated audit CSV: \path{results/supplementary_evidence/decision_aware_robustness_summary.csv}.
\endgroup
\end{table}

% Supplementary table 4: tab:threats_validity_matrix
\begin{table}[t]
\centering
\caption{Threats-to-Validity matrix for reviewer-facing closure.}
\label{tab:threats_validity_matrix}
\begingroup
\setlength{\tabcolsep}{3pt}
\scriptsize
\resizebox{\linewidth}{!}{%
\begin{tabular}{L{0.17\linewidth}L{0.21\linewidth}L{0.20\linewidth}L{0.20\linewidth}L{0.20\linewidth}}
\toprule
Reviewer concern & Current evidence & Residual limitation & Claim boundary & Planned future closure \\
\midrule
Model non-dominance & Main matrix + dispatch matrix show CEEMDAN-WOA-GRU/iTransformer wins on selected metrics. & Ultra-LSNT is not best on clean or mapped dispatch ranking. & Bounded-compute robustness operating point, not universal best. & Expand multi-dataset decision-focused ranking with explicit utility weights. \\
Mapping dependence & Raw-vs-mapped gate + three-point aggregates (\path{dispatch_mapping_sensitivity_wmax260_280_300.csv}) with day-level consistency evidence (\path{dispatch_mapping_daylevel_consistency_wmax260_280_300.csv}). & Results depend on fixed mapping protocol and are not natural system-readiness proof. & Dispatch conclusions are mapped-protocol conditional evidence. & Add broader mapping-parameter sweep and site-specific calibration appendix. \\
Fairness boundary & Symmetric clean/fault-augmented closure for four principal models (DLinear, TimeMixer, iTransformer, Ultra-LSNT; 3 seeds; 4 faults) from \path{fairness_closure_fullfault_hard_table.csv}. & Long-tail supplementary baselines remain clean-trained or constrained-budget diagnostics. & Main comparative claims are restricted to the four symmetry-closed models; long-tail baselines are supplementary diagnostics only. & Extend full symmetry-closure protocol to all long-tail baselines. \\
Physics boundary & Three-way guardrail control (\path{fgmoe_clip_only_control_windcn.csv}) plus FG-MoE boundary diagnostics. & Guardrail is heuristic envelope control, not full mechanism model. & Feasibility-guardrail positioning, not PDE-constrained solver claim. & Add lightweight aerodynamic surrogate and stronger physics coupling. \\
Hardware closure & CPU/ARM screening with direct + backfilled split (\path{latency_arm_full_backfill.csv}). & Not fully normalized bake-off; no board-level HIL closure yet. & Cross-ISA deployability precondition evidence only. & Add PLC/RTU board-level HIL latency, thermal, and power profiling. \\
AC/system-grade closure & RTS-24 DC-OPF screening aggregates (\path{dispatch_network_ieee24_dcopf_wmax260_aggregate.csv}, \path{dispatch_network_ieee24_dcopf_wmax280_aggregate.csv}, \path{dispatch_network_ieee24_dcopf_wmax300_aggregate.csv}). & DC layer is comparative proxy, not AC security validation. & No operator-ready AC-feasibility certification claim. & Joint SCUC/AC-OPF + voltage-stability diagnostics. \\
\bottomrule
\end{tabular}% 
}
\vspace{3pt}
\footnotesize Evidence-path index used above: fairness (main-claim set) \path{results/supplementary_evidence/fairness_closure_fullfault_hard_table.csv}, diagnostic extension \path{results/supplementary_evidence/fairness_minimal_aug_controls_drift_2seed.summary.csv}; mapping aggregates \path{results/supplementary_evidence/dispatch_mapping_sensitivity_wmax260_280_300.csv}; mapping day-level consistency \path{results/supplementary_evidence/dispatch_mapping_daylevel_consistency_wmax260_280_300.csv}; physics guardrail \path{results/supplementary_evidence/fgmoe_clip_only_control_windcn.csv}; hardware backfill \path{results/supplementary_evidence/latency_arm_full_backfill.csv}; dispatch network \path{results/supplementary_evidence/dispatch_network_ieee24_dcopf_wmax260_aggregate.csv}, \path{results/supplementary_evidence/dispatch_network_ieee24_dcopf_wmax280_aggregate.csv}, and \path{results/supplementary_evidence/dispatch_network_ieee24_dcopf_wmax300_aggregate.csv}.
\endgroup
\end{table}

% Supplementary table 3: tab:pimoe_windcn_ablation
\begin{table}[t]
\centering
\caption{Ultra-LSNT and FG-MoE comparison on Wind (CN) under clean input and severe Gaussian corruption ($\sigma_{\mathrm{eff}}=0.6$). RMSE/MAE are in kW after inverse scaling.}
\label{tab:pimoe_windcn_ablation}
\begin{tabular}{llccc}
\toprule
Setting & Model & $R^2$ & RMSE & MAE \\
\midrule
Clean & Ultra-LSNT & 0.8727 & 9081.40 & 7168.69 \\
Clean & FG-MoE (full-trained) & \textbf{0.8747} & \textbf{9012.34} & \textbf{6986.37} \\
Gaussian noise ($\sigma_{\mathrm{eff}}=0.6$) & Ultra-LSNT & 0.8488 & 9900.07 & 7871.25 \\
Gaussian noise ($\sigma_{\mathrm{eff}}=0.6$) & FG-MoE (full-trained) & \textbf{0.8526} & \textbf{9774.56} & \textbf{7691.53} \\
\bottomrule
\end{tabular}
\end{table}

% Supplementary table 4: tab:pimoe_guardrail_metrics
\begin{table}[t]
\centering
\caption{Physical-violation diagnostics for Ultra-LSNT and FG-MoE on Wind (CN); percentages are computed on inverse-scaled predictions under a clip-control protocol (supplementary mechanism diagnostic). Here the deployable proxy is $v_{t+h}^{\mathrm{proxy}}=v_t$ and rated-power exceedance means $\hat y>P_{\mathrm{rated}}$.}
\label{tab:pimoe_guardrail_metrics}
\begingroup
\setlength{\tabcolsep}{4pt}
\small
\resizebox{\linewidth}{!}{%
\begin{tabular}{llccc}
\toprule
Setting & Model & Negative (\%) & Rated-power exceed. (\%) & Cut-region nonzero (\%) \\
\midrule
Clean & Ultra-LSNT & 0.0000 & 0.1164 & 0.9387 \\
Clean & FG-MoE (full-trained) & 0.0000 & \textbf{0.0000} & 0.9387 \\
$\sigma_{\mathrm{eff}}=0.6$ & Ultra-LSNT & 0.0000 & 0.0653 & \textbf{1.6896} \\
$\sigma_{\mathrm{eff}}=0.6$ & FG-MoE (full-trained) & 0.0000 & \textbf{0.0023} & 1.7558 \\
\bottomrule
\end{tabular}%
}
\vspace{3pt}
\footnotesize Protocol note: this supplementary table uses the clip-control evidence path (\path{fgmoe_clip_only_control_windcn.csv}) with the deployable proxy $v_{t+h}^{\mathrm{proxy}}=v_t$, and is not directly comparable to the main-text safety values in Table~4 Panel B, which use protocol P1 (\path{physics_closure.csv}).
\endgroup
\end{table}

% Supplementary table 5: tab:robust_gaussian_sweep
\begin{table}[t]
\centering
\caption{Gaussian-noise robustness sweep on Wind (CN) with configured ratio $\sigma_{\mathrm{cfg}}$ and effective ratio $\sigma_{\mathrm{eff}}=\alpha\sigma_{\mathrm{cfg}}$ ($\alpha=1.5$).}
\label{tab:robust_gaussian_sweep}
\begingroup
\setlength{\tabcolsep}{3pt}
\scriptsize
\resizebox{\linewidth}{!}{%
\begin{tabular}{lccccccc}
\toprule
$\sigma_{\mathrm{cfg}}$ ($\sigma_{\mathrm{eff}}$) & Ultra-LSNT $R^2$ & iTransformer $R^2$ & TimeMixer $R^2$ & PatchTST $R^2$ & DLinear $R^2$ & LightGBM $R^2$ & Persistence $R^2$ \\
\midrule
0.00 (0.00) & 0.8728 & 0.9012 & 0.8886 & 0.8666 & \textbf{0.9379} & 0.9069 & 0.9175 \\
0.10 (0.15) & 0.8687 & 0.8993 & 0.8847 & 0.8659 & 0.7814 & 0.8860 & 0.9140 \\
0.20 (0.30) & 0.8552 & 0.8937 & 0.8738 & 0.8632 & 0.2952 & 0.8222 & 0.9034 \\
0.30 (0.45) & 0.7866 & 0.8844 & 0.8566 & 0.8582 & -0.4659 & 0.7305 & 0.8879 \\
0.40 (0.60) & 0.7263 & 0.8717 & 0.8337 & 0.8508 & -1.5790 & 0.6156 & 0.8612 \\
\bottomrule
\end{tabular}%
}
\endgroup
\end{table}

% Supplementary table 6: tab:fairness_drift_2seed
\begin{table}[t]
\centering
\caption{Minimal matched robustness controls for three strong baselines using one non-Gaussian perturbation (drift) and two seeds. This is supplementary diagnostic evidence only and is not used for main comparative claims, which rely on the four-model full-fault closure matrix.}
\label{tab:fairness_drift_2seed}
\label{tab:itransformer_augmented}
\begin{tabular}{lcccc}
\toprule
Model & Clean baseline $R^2$ (drift@0.60) & Drift augmented $R^2$ (drift@0.60) & Gain & $n_{\text{seed}}$ \\
\midrule
DLinear & $0.8730 \pm 0.0000$ & $\mathbf{0.8999 \pm 0.0004}$ & $+0.0269$ & 2 \\
iTransformer & $0.8267 \pm 0.0000$ & $\mathbf{0.8773 \pm 0.0090}$ & $+0.0505$ & 2 \\
TimeMixer & $0.8263 \pm 0.0000$ & $\mathbf{0.8622 \pm 0.0028}$ & $+0.0359$ & 2 \\
\bottomrule
\end{tabular}
\vspace{3pt}
\footnotesize Source: \path{results/supplementary_evidence/fairness_minimal_aug_controls_drift_2seed.summary.csv}. Claim boundary: diagnostic extension only; main claims use \path{results/supplementary_evidence/fairness_closure_fullfault_hard_table.csv}.
\end{table}

% Supplementary table 7: tab:eff_robust_tradeoff
\begin{table}[t]
\centering
\caption{Guardrail mechanism control on Wind (CN): no-guardrail, posthoc clip-only baseline, and FG-MoE (full-trained). Clean and severe Gaussian ($\sigma_{\mathrm{eff}}=0.60$) slices are reported under the clip-control supplementary protocol with $v_{t+h}^{\mathrm{proxy}}=v_t$. The clip-only row is a posthoc screening baseline, not the FG-MoE main method, and should not be directly compared with main-text Table~4 Panel B (protocol P1).}
\label{tab:clip_only_guardrail_control}
\begingroup
\setlength{\tabcolsep}{4pt}
\small
\begin{tabular}{llcccc}
\toprule
Scenario & Model & $R^2$ & RMSE (kW) & MAE (kW) & Rated-power exceed. (\%) \\
\midrule
Clean & Ultra-LSNT (no guardrail) & 0.8727 & 9081.40 & 7168.69 & 0.1164 \\
Clean & Ultra-LSNT + clip-only (posthoc) & 0.4319 & 19189.01 & 14757.99 & 0.0000 \\
Clean & FG-MoE (full-trained) & 0.8747 & 9012.34 & 6986.37 & 0.0000 \\
\midrule
$\sigma_{\mathrm{eff}}=0.60$ & Ultra-LSNT (no guardrail) & 0.8502 & 9852.04 & 7849.67 & 0.0656 \\
$\sigma_{\mathrm{eff}}=0.60$ & Ultra-LSNT + clip-only (posthoc) & 0.3773 & 20088.82 & 15335.83 & 0.0000 \\
$\sigma_{\mathrm{eff}}=0.60$ & FG-MoE (full-trained) & 0.8526 & 9774.56 & 7691.53 & 0.0023 \\
\bottomrule
\end{tabular}
\vspace{3pt}
\footnotesize Source: \path{results/supplementary_evidence/fgmoe_clip_only_control_windcn.csv}; FG-MoE rows are aligned to \path{results/supplementary_evidence/pi_moe_full_vs_ultra_clean_sigma06_windcn.csv}. Rated-power exceedance is defined as $\hat y>P_{\mathrm{rated}}$, and the shared deployable proxy is $v_{t+h}^{\mathrm{proxy}}=v_t$. Protocol boundary: this table is supplementary mechanism diagnostics and is not a direct numerical counterpart of Table~4 Panel B (\path{physics_closure.csv}).
\endgroup
\end{table}

% Supplementary table 7: tab:eff_robust_tradeoff
\begin{table}[t]
\centering
\caption{Compact efficiency-robustness summary on Wind (CN) using clean and noisy ($\sigma_{\mathrm{eff}}=0.30$) settings.}
\label{tab:eff_robust_tradeoff}
\scriptsize
\begin{tabular}{lcccc}
\toprule
Model & Latency (ms) & Clean $R^2$ & Noisy $R^2$ & $R^2$ drop (\%) \\
\midrule
DLinear & 0.265 & 0.9379 & 0.2952 & 68.5 \\
LightGBM & 50.397 & 0.9069 & 0.8222 & 9.3 \\
Ultra-LSNT & 3.714 & 0.8728 & 0.8552 & 2.0 \\
\bottomrule
\end{tabular}
\end{table}

% Supplementary table 8: tab:robust_structured_sev06
\begin{table}[t]
\centering
\caption{Robustness under structured SCADA faults on Wind (CN) at severity 0.6 ($R^2$ after inverse scaling).}
\label{tab:robust_structured_sev06}
\begin{tabular}{lccc}
\toprule
Fault & Ultra-LSNT $R^2$ & iTransformer $R^2$ & DLinear $R^2$ \\
\midrule
Flatline & 0.8665 & 0.8942 & \textbf{0.9199} \\
Spike & 0.7829 & \textbf{0.8609} & 0.2058 \\
Dropout & 0.8667 & 0.8947 & \textbf{0.9245} \\
\bottomrule
\end{tabular}
\end{table}

% Supplementary table 9: tab:robust_missingness
\begin{table}[t]
\centering
\caption{Robustness under missingness and time-axis imputation on Wind (CN) (3 seeds, mean$\pm$std).}
\label{tab:robust_missingness}
\begingroup
\setlength{\tabcolsep}{4pt}
\small
\resizebox{\linewidth}{!}{%
\begin{tabular}{lccccc}
\toprule
Imputation & $p$ & Ultra-LSNT $R^2$ & Ultra-LSNT RMSE (kW) & DLinear $R^2$ & DLinear RMSE (kW) \\
\midrule
\texttt{ffill} & 0.1 & $0.8718\pm0.0001$ & $9115.0\pm3.6$ & $0.9326\pm0.0001$ & $6610.9\pm4.4$ \\
\texttt{ffill} & 0.3 & $0.8692\pm0.0002$ & $9205.8\pm8.4$ & $0.9242\pm0.0005$ & $7010.0\pm23.9$ \\
\texttt{linear} & 0.1 & $0.8722\pm0.0001$ & $9102.4\pm2.9$ & $0.9325\pm0.0001$ & $6613.6\pm6.0$ \\
\texttt{linear} & 0.3 & $0.8705\pm0.0002$ & $9162.9\pm6.9$ & $0.9239\pm0.0006$ & $7021.1\pm27.9$ \\
\bottomrule
\end{tabular}%
}
\endgroup
\end{table}

% Supplementary table 10: tab:multi_domain_baselines
\begin{table}[t]
\centering
\caption{Dataset-native error metrics for Ultra-LSNT after inverse scaling; units follow each dataset.}
\label{tab:multi_domain_baselines}
\begin{tabular}{lcccc}
\toprule
Dataset & RMSE & MAE & Unit & Skip rate (\%)\\
\midrule
Wind (CN) & 9081 & 7171 & kW & 0.0 \\
Wind (US) & 0.4556 & 0.2891 & MW & 11.4 \\
Air Quality & 32.21 & 19.87 & AQI & 0.0 \\
GEFCom Load & 2837 & 2049 & MW & 0.0 \\
\bottomrule
\end{tabular}
\end{table}

% Supplementary table 11: tab:ramp_tail_smoothness
\begin{table}[t]
\centering
\caption{Ramp, tail-error, smoothness, and routing-stability diagnostics on Wind (CN). Values are reported for auditable Ultra-LSNT and DLinear runs.}
\label{tab:ramp_tail_smoothness}
\begingroup
\setlength{\tabcolsep}{3.5pt}
\small
\resizebox{\linewidth}{!}{%
\begin{tabular}{lrrrrrrr}
\toprule
Model & $P95(|e|)$ (kW) & $\max(|e|)$ (kW) & Ramp F1 & TV (kW) & TV/TV$_{\text{true}}$ & Routing var & (max-share, entropy) \\
\midrule
Ultra-LSNT & 11,940.17 & 34,544.14 & 0.3544 & 8,222,569.29 & 0.2525 & 0.04413 & (0.5394, 0.7511) \\
DLinear & 2,484.47 & 21,479.02 & 0.5939 & 29,013,519.33 & 0.8909 & 0.00000 & (1.0000, 0.0000) \\
\bottomrule
\end{tabular}%
}
\endgroup
\end{table}

% Supplementary table 12: tab:prob_gaussian_minimal
\begin{table}[t]
\centering
\caption{Post-hoc Gaussian calibration baseline on Wind (CN) with nominal 80\% and 90\% intervals.}
\label{tab:prob_gaussian_minimal}
\begingroup
\setlength{\tabcolsep}{4pt}
\small
\resizebox{\linewidth}{!}{%
\begin{tabular}{lrrrrrr}
\toprule
Model & $\sigma$ (kW) & NLL & CRPS & PICP@80 / PINAW@80 & PICP@90 / PINAW@90 \\
\midrule
Ultra-LSNT & 5,044.08 & 10.1314 & 3,353.4667 & 0.7253 / 0.10995 & 0.8242 / 0.14111 \\
DLinear & 1,329.74 & 8.6118 & 670.2737 & 0.8726 / 0.02898 & 0.9269 / 0.03720 \\
\bottomrule
\end{tabular}%
}
\endgroup
\end{table}

% Supplementary table 13: tab:cpu_batch1_latency
\begin{table}[t]
\centering
\caption{Single-thread CPU batch-1 microbenchmark for online inference latency.}
\label{tab:cpu_batch1_latency}
\begin{tabular}{lcc}
\toprule
Model & Latency (ms) & Throughput (Hz) \\
\midrule
Ultra-LSNT (full) & 3.714 $\pm$ 0.341 & 269.3 \\
Ultra-LSNT-Lite (orig.) & 2.430 $\pm$ 0.245 & 411.5 \\
Ultra-LSNT-Lite (opt.) & 1.171 $\pm$ 0.201 & 854.0 \\
VMD-PSO-LSTM (minimal auditable) & 10.297 $\pm$ 0.791 & 97.1 \\
CEEMDAN-WOA-GRU (minimal auditable) & 11.082 $\pm$ 0.372 & 90.2 \\
DLinear & 0.2654 $\pm$ 0.0843 & 3767.6 \\
LightGBM & 50.3970 $\pm$ 21.1347 & 19.8 \\
\bottomrule
\end{tabular}
\end{table}

% Supplementary table 14: tab:arm_microbenchmark
\begin{table}[t]
\centering
\caption{ARMv8 (aarch64) batch-1 microbenchmark on Huawei Kunpeng~920 for cross-ISA latency screening (archived logs).}
\label{tab:arm_microbenchmark}
\begin{tabular}{lcc}
\toprule
Model & Avg latency (ms) & P95 latency (ms) \\
\midrule
Ultra-LSNT (Full) & 8.795 & 10.419 \\
Ultra-LSNT-Lite (orig.) & 4.632 & 5.390 \\
Ultra-LSNT-Lite (opt.) & 5.007 & 5.900 \\
DLinear & 0.102 & 0.105 \\
TimeMixer & 1.811 & 2.120 \\
iTransformer & 1.382 & 1.943 \\
\bottomrule
\end{tabular}
\end{table}

% Supplementary table 15: tab:tau_sweep_pareto
\begin{table}[t]
\centering
\caption{Jump Gate threshold sweep on Wind (CN) with $R^2$, RMSE, skip rate, and CPU batch-1 latency.}
\label{tab:tau_sweep_pareto}
\begin{tabular}{lcccc}
\toprule
$\tau$ & $R^2$ & RMSE (kW) & Skip rate (\%) & CPU latency (ms) \\
\midrule
0.3 (Edge-Frugal) & 0.8661 & 9314.40 & 33.38 & 3.621 \\
0.5 (Reliability, default) & 0.8728 & 9081.66 & 0.00 & 3.628 \\
0.7 & 0.8728 & 9081.66 & 0.00 & 3.766 \\
0.9 & 0.8728 & 9081.66 & 0.00 & 3.665 \\
\bottomrule
\end{tabular}
\end{table}

% Supplementary table 16: tab:impl_hardware
\begin{table}[t]
\centering
\caption{Computational hardware used for training and inference profiling.}
\label{tab:impl_hardware}
\begin{tabular}{ll}
\toprule
Hardware component & Specification \\
\midrule
GPU & NVIDIA GeForce RTX 4090 (or equivalent) \\
CPU & Intel Core i9-13900K (or equivalent) \\
Memory & 64\,GB DDR5 RAM \\
Storage & 2\,TB NVMe SSD \\
OS & Ubuntu 22.04 LTS / Windows 11 (WSL2) \\
\bottomrule
\end{tabular}
\end{table}

% Supplementary table 17: tab:model_config
\begin{table}[t]
\centering
\caption{Ultra-LSNT and Ultra-LSNT-Lite configuration summary used for reproducibility.}
\label{tab:model_config}
\begingroup
\setlength{\tabcolsep}{4pt}
\small
\resizebox{\linewidth}{!}{%
\begin{tabular}{lcc}
\toprule
Hyperparameter & Ultra-LSNT (main) & Ultra-LSNT-Lite (Opt.) \\
\midrule
Lookback $L$ / horizon $H$ & 96 / 24 & 96 / 24 \\
Model dimension $d$ & 128 & 96 \\
\# MoE blocks / layers & 3 & 1 \\
\# experts $N$ & 4 & 2 \\
Top-$K$ routing $K$ & 2 & 1 \\
Router auxiliary loss $\lambda$ & 0.01 & 0.01 \\
Jump Gate threshold $\tau$ & 0.5 & 0.5 \\
Dropout & 0.1 & 0.1 \\
Optimizer / LR & AdamW / $1\times10^{-3}$ & AdamW / $1\times10^{-3}$ \\
Batch size & 128 & 128 \\
Inference optimization & Eager FP32 (no compile/AMP) & \texttt{torch.compile}+AMP (FP16) \\
\bottomrule
\end{tabular}%
}
\endgroup
\end{table}

% Supplementary table 18: tab:robust_structured_full
\begin{table}[t]
\centering
\caption{Full structured SCADA-fault robustness results on Wind (CN) for Ultra-LSNT, iTransformer, and DLinear.}
\label{tab:robust_structured_full}
\begingroup
\setlength{\tabcolsep}{3pt}
\scriptsize
\resizebox{\linewidth}{!}{%
\begin{tabular}{lllccccc}
\toprule
Fault type & Severity & Model & $R^2$ & RMSE (kW) & MAE (kW) & $\Delta R^2_{\text{abs}}$ & $\Delta R^2_{\text{rel}}$ (\%) \\
\midrule
clean & 0.0 & Ultra-LSNT & 0.8727 & 9081.4 & 7168.7 & 0.0000 & 0.00 \\
clean & 0.0 & iTransformer & 0.9012 & 8004.0 & 6333.3 & 0.0000 & 0.00 \\
clean & 0.0 & DLinear & 0.9379 & 6342.9 & 4203.5 & 0.0000 & 0.00 \\
\midrule
flatline & 0.2 & Ultra-LSNT & 0.8721 & 9103.3 & 7181.3 & 0.0006 & 0.07 \\
flatline & 0.2 & iTransformer & 0.9005 & 8029.5 & 6351.9 & 0.0006 & 0.07 \\
flatline & 0.2 & DLinear & 0.9369 & 6394.6 & 4230.8 & 0.0010 & 0.11 \\
\midrule
flatline & 0.6 & Ultra-LSNT & 0.8665 & 9302.9 & 7345.3 & 0.0063 & 0.72 \\
flatline & 0.6 & iTransformer & 0.8942 & 8278.8 & 6541.4 & 0.0069 & 0.77 \\
flatline & 0.6 & DLinear & 0.9199 & 7203.7 & 4639.0 & 0.0180 & 1.92 \\
\midrule
spike & 0.2 & Ultra-LSNT & 0.8509 & 9831.3 & 7737.6 & 0.0219 & 2.51 \\
spike & 0.2 & iTransformer & 0.8897 & 8454.4 & 6594.6 & 0.0114 & 1.27 \\
spike & 0.2 & DLinear & 0.7766 & 12031.7 & 5335.3 & 0.1613 & 17.20 \\
\midrule
spike & 0.6 & Ultra-LSNT & 0.7829 & 11860.9 & 9324.8 & 0.0898 & 10.29 \\
spike & 0.6 & iTransformer & 0.8609 & 9494.3 & 7228.6 & 0.0402 & 4.46 \\
spike & 0.6 & DLinear & 0.2058 & 22686.9 & 7864.8 & 0.7321 & 78.05 \\
\midrule
dropout & 0.2 & Ultra-LSNT & 0.8721 & 9105.5 & 7190.2 & 0.0007 & 0.08 \\
dropout & 0.2 & iTransformer & 0.9004 & 8034.7 & 6363.1 & 0.0008 & 0.08 \\
dropout & 0.2 & DLinear & 0.9358 & 6450.5 & 4275.6 & 0.0021 & 0.23 \\
\midrule
dropout & 0.6 & Ultra-LSNT & 0.8667 & 9294.1 & 7334.4 & 0.0060 & 0.69 \\
dropout & 0.6 & iTransformer & 0.8947 & 8259.4 & 6537.0 & 0.0064 & 0.71 \\
dropout & 0.6 & DLinear & 0.9245 & 6993.4 & 4557.9 & 0.0134 & 1.43 \\
\bottomrule
\end{tabular}%
}
\endgroup
\end{table}

% Supplementary table 19: tab:ablation_kernel
\begin{table}[t]
\centering
\caption{Kernel-size ablation on Wind (CN) from archived ablation runs.}
\label{tab:ablation_kernel}
\begin{tabular}{lc}
\toprule
Kernel size & $R^2$ \\
\midrule
3 & 0.8367 \\
7 & 0.8360 \\
15 & 0.8236 \\
\bottomrule
\end{tabular}
\end{table}

% Supplementary table 20: tab:expert_physics
This interpretability block supports the main-text claim boundary that routing evidence is used for mechanism auditability rather than for additional accuracy-leaderboard claims; readers should consult Table~\ref{tab:expert_physics} to verify expert-physics alignment.
\begin{table}[t]
\centering
\caption{Expert specialization summary using the highest-correlated physical feature.}
\label{tab:expert_physics}
\begin{tabular}{lcc}
\toprule
Expert & Feature & Correlation \\
\midrule
Expert 1 & wind\_speed & +0.42 \\
Expert 2 & turbulence & +0.38 \\
Expert 3 & temperature & +0.31 \\
Expert 4 & humidity & +0.25 \\
\bottomrule
\end{tabular}
\end{table}

% Supplementary table 21: tab:dm_test
This significance block supports the main-text stability argument by testing paired loss-direction consistency; readers should use Table~\ref{tab:dm_test} to check whether direction-level differences are statistically reliable under fixed splits.
\begin{table}[t]
\centering
\caption{Diebold--Mariano test results on Wind (CN) using squared-error loss; DM$<0$ favors Ultra-LSNT.}
\label{tab:dm_test}
\begin{tabular}{lcc}
\toprule
Comparison & DM statistic & $p$-value \\
\midrule
Ultra-LSNT vs PatchTST & -5.90 & $1.82 \times 10^{-9}$ \\
Ultra-LSNT vs TimeMixer & 3.45 & $2.83 \times 10^{-4}$ \\
Ultra-LSNT vs iTransformer & 5.12 & $3.07 \times 10^{-7}$ \\
\bottomrule
\end{tabular}
\end{table}

\end{document}

